\subsection{\texorpdfstring{有界测度与 \(\mathcal{C}^{b}(T)\) 上的线性形式}{有界测度与 \textbackslash mathcal\{C\}\^{}\{b\}(T) 上的线性形式}}

命题 5. --- 设 T 是完全正则空间，I 是赋范空间 \(\mathcal{C}^{b}(T;C)\) 上的连续复线性形式。要使 T 上存在有界复测度 \(\theta\)，满足 \(\theta(f)=\mathrm{I}(f)\) 对所有 \(f\in\mathcal{C}^{b}(T;C)\) 成立，以下条件成立是必要且充分的：

\begin{enumerate}
\def\labelenumi{(\Alph{enumi})}
\setcounter{enumi}{12}
\tightlist
\item
  对每个数 \(\varepsilon > 0\)，存在紧子集 \(K\)（属于 \(T\)），使得关系 \(g \in \mathcal{C}^b(T; \mathbf{C})\)、\(|g| \leqslant 1\)、\(g_K = 0\) 蕴含 \(|\mathrm{I}(g)| \leqslant \varepsilon\)。
\end{enumerate}

测度 \(\theta\) 于是唯一。

唯一性由第 1 号命题 2 得出。来证明条件 (M) 的必要性。设 \(\theta\) 是有界复测度；设 K 是满足 \(|\theta|^{\bullet}(\mathrm{T}-\mathrm{K}) \leqslant \varepsilon\) 的紧集（§1, No.~2, 备注 3）。假设 \(|g| \leqslant 1\)、\(g_{\mathbf{K}} = 0\)，则有 \(|g| \leqslant \varphi_{\mathbf{C}_{\mathbf{K}}}\)，因而 \(|\theta(g)| \leqslant |\theta|^{\bullet}(\varphi_{\mathbf{C}_{\mathbf{K}}}) \leqslant \varepsilon\)。

现在证明充分性。令 X 为 T 的 Stone--Čech 紧化（GT, IX, §1, 习题 7；或 TG, IX, §1, No.~6）。对每个函数 \(f \in \mathcal{C}(X; \mathbf{C})\)，置 \(\nu(f) = I(f_{\mathrm{T}})\)；这样就在 \(\mathcal{C}(X; \mathbf{C})\) 上定义了一个连续线性形式，也就是 X 这个紧空间上的复测度。取 \(\varepsilon\) 为满足 \(>0\) 的数，K 为满足 (M) 的紧集；由于函数 \(\varphi_{\mathbf{C}_{\mathrm{K}}}\) 在 X 上下半连续且为正，公式 (2) 给出以下关系，其中 \(\mathcal{G}\) 表示由满足 \(g \in \mathcal{C}(X; \mathbf{C})\) 的函数 \(|g| \leqslant \varphi_{\mathbf{C}_{\mathrm{K}}}\) 构成的集合：

\[
| \nu | ^ {\bullet} (\mathrm{X} - \mathrm{K}) = \sup _ {g \in \mathcal {G}} | \nu (g) | = \sup _ {g \in \mathcal {G}} | \mathrm{I} (g _ {\mathrm{T}}) | \leqslant \varepsilon .
\]

设 \((\mathrm{K}_n)_{n\geqslant 1}\) 是 T 的紧子集序列，使得每个 \(\mathbf{K}_n\) 对 \(\varepsilon = 1 / n\) 满足 (M)，并令 \(\mathrm{S} = \bigcup_{n}\mathrm{K}_{n}\)；S 是 X 中包含于 T 的 Borel 集，并且 \(|\nu|^{\bullet}(\mathrm{X} - \mathrm{T})\leqslant |\nu |^{\bullet}(\mathrm{X} - \mathrm{S})\leqslant |\nu |^{\bullet}(\mathrm{X} - \mathrm{K}_n)\leqslant 1 / n\) 对所有 \(n\) 成立，所以 T 是 \(\nu\)-可测的，且 \(\nu\) 集中于 T。设 \(f\) 是 T 上的有界连续函数；由于 X 是 T 的 Stone-Čech 紧化，\(f\) 可连续延拓为函数 \(g\in \mathcal{C}(\mathrm{X};\mathbf{C})\)。现在令 \(\mu\) 为 \(\nu\) 在 T 上诱导的测度；有 \(\mu(f) = \nu(f^0)\)。由于 \(\nu\) 集中于 T，函数 \(f^0\) 与 \(g\) 几乎处处相等（相对于 \(\nu\)），因此 \(\mu(f) = \nu(g) = I(g_T) = I(f)\)，证明完毕。

推论。 --- 沿用命题 5 的记号，设 T 上存在有界正测度 \(\mu\)，使得 \(|\mathrm{I}(f)| \leqslant \mu(|f|)\) 对所有 \(f \in \mathcal{C}^b(\mathrm{T};\mathbf{C})\) 成立；则 T 上存在复测度 \(\theta\)，使得 \(\theta(f) = \mathrm{I}(f)\) 对所有 \(f \in \mathcal{C}^b(\mathrm{T};\mathbf{C})\) 成立。

